%============================= abstract.tex================================
\chapter*{Abstract}%
%\addcontentsline{toc}{chapter}{\numberline{}Abstract}%
\addcontentsline{toc}{chapter}{Abstract}%

Legacy scientific and engineering software is predominantly written as scalar loop nests in C, C++, and Fortran, which makes direct use of modern tensor accelerators difficult without expensive manual rewriting. This project implements LoopHole, a verification-driven lifting tool that translates loop-level MLIR Affine programs into higher-level MLIR tensor dialects, primarily Linalg and partially StableHLO. The implemented pipeline consists of deterministic Affine extraction, sketch-based candidate generation, SymPy-guided confidence scoring, formal equivalence checking using Z3, and dialect-specific code emission through a command-line interface. The current implementation supports practical kernels such as matrix multiplication, transpose, dot product, 1D and 2D convolution, elementwise addition, and reduction patterns. In a strict-profile batch run on seven canonical fixture programs, the system produced seven accepted outputs with formal equivalence proofs and generated lifted MLIR artifacts for all processed inputs. The recorded verification times ranged from 5.27 ms to 33.98 ms, with an average of approximately 12.87 ms, indicating that the current proof strategy is efficient for constrained dense kernels. The project also provides reproducible workflows through Dockerized Polygeist and structured batch reports for weekly benchmarking and regression tracking. Overall, this work demonstrates the practical feasibility of automated lifting with verification-aware synthesis in the MLIR ecosystem, while also identifying open challenges in StableHLO parity, dynamic-shape proof strategies, and broader real-world kernel coverage.
\thispagestyle{plain}
%=======================================================================
